% Course material: formal methods with the Z notation.
% Deliberately carries a metadata/changelog preamble so the pollution
% diagnostic has realistic structural distractors to classify.

\documentclass{article}
\title{Predicate Logic in Z}
\author{Formal Methods Course}
\date{Revision history below}

% Changelog (structural metadata, not substantive content):
%   2026-01-04  Initial draft of the predicate-logic chapter.
%   2026-02-11  Added worked examples for the one-point rule.
%   2026-03-02  Fixed a typo in the universal-quantifier schema.

\begin{document}
\maketitle

\section{Introduction}
The Z notation builds specifications on typed set theory and first-order
predicate logic. Before schemas and state, a specifier needs the predicate
calculus: quantifiers, bound and free variables, and the laws that let one
predicate be rewritten as another while preserving meaning.

\section{Predicate Logic in Z}
A predicate in Z is a boolean-valued expression over typed variables. The
universal quantifier is written $\forall x : T \bullet P$, read ``for all
$x$ of type $T$, the predicate $P$ holds.'' The existential quantifier
$\exists x : T \bullet P$ asserts that some $x$ of type $T$ makes $P$ true.
Quantifiers introduce a declaration, a constraint, and a predicate body,
separated by a vertical bar and a bullet: $\forall x : T | C \bullet P$.

A variable is bound when a quantifier ranges over it and free otherwise.
The one-point rule eliminates an existential quantifier when the bound
variable is fixed by an equality: $\exists x : T \bullet x = e \land P$
simplifies to $P[e/x]$ provided $e$ has type $T$. De Morgan's laws relate
the quantifiers: $\lnot (\forall x : T \bullet P)$ equals
$\exists x : T \bullet \lnot P$.

\section{Schemas and State}
A schema packages a declaration part and a predicate part. The predicate
part is ordinary predicate logic over the declared variables, so every
law from the previous section applies unchanged inside a schema.

\section{Sequences and Bags}
Sequences model ordered collections and bags model multisets. Both are
defined as functions, and reasoning about them again reduces to the
predicate calculus over their defining functions.

\end{document}
